% MIT; illustrative data. Compile from repository root.
\documentclass[UTF8,fontset=fandol]{ctexart}
\usepackage[margin=18mm]{geometry}
\usepackage{gradbars}
\setlength{\parindent}{0pt}
\pagestyle{empty}
\begin{document}
\section*{区域项目评审}
\gradbarsloadcsv{projects}{docs/data/hierarchical-teams.csv}
\begin{center}
\begin{gradtable}[data=projects,style=hierarchical,
 width=\linewidth,group levels={region,team},
 rank scope=parent,parent summary=true,leaf summary=true,
 parent summary label={区域均值},
 leaf summary label={小组均值},label width=17mm]
\gradcolumn[source=Project,width=22mm]{project}{项目}
\gradcolumn[source=Region,width=17mm]{region}{区域}
\gradcolumn[source=Team,width=18mm]{team}{小组}
\gradcolumn[source=Score,type=number,width=20mm,
 aggregate=mean,mark=both,options={precision=1}]{score}{得分}
\gradcolumn[source=Trend,type=trendcompare,
 options={min=0,max=100,height=5mm,precision=0,
 compare label=delta}]{trend}{首末 / 趋势}
\gradsummary{可见行均值}
\end{gradtable}
\end{center}
\bigskip
\begingroup
\gradscale{project-score}{min=0,max=100}
\gradbarssetup{width=65mm,height=5ex,
 precision=0,label width=20mm,compare label=delta}
\begin{tabular}{lc}
项目 & 首末对比与全过程趋势\\
\hline
Atlas & \gradsparkcompare[scale group=project-score]
 {20,40,NA,80}\\
Birch & \gradsparkcompare[scale group=project-score]
 {NA,30,60,90}\\
Elm & \gradsparkcompare[scale group=project-score]
 {20,50,65,NA}\\
\end{tabular}
\endgroup
\par\medskip
不同区域下的同名小组保持独立；首末缺失不回退到中间有效值。
\end{document}
